\ExerciseSection

\begin{enumerate}
\def\labelenumi{\arabic{enumi})}
\tightlist
\item
  Let \(f\) be the mapping of \(\mathbf{S}_2\) into \(\mathbb{R}^4\) such that
\end{enumerate}

\[
f \left(x _ {1}, x _ {2}, x _ {3}\right) = \left(x _ {1} ^ {2} - x _ {2} ^ {2}, x _ {1} x _ {2}, x _ {1} x _ {3}, x _ {2} x _ {3}\right).
\]

This function has the same value at diametrically opposite points of \(S_{2}\) . Show that on passing to the quotient it defines a homeomorphism of \(P_{2}\) onto a subspace of \(R^{4}\) .

Show likewise that the mapping \(g\) of \(\mathbf{S}_3\) into \(\mathbf{R}^6\) defined by

\[
g \left(x _ {1}, x _ {2}, x _ {3}, x _ {4}\right) = \left(x _ {1} ^ {2} - x _ {2} ^ {2}, x _ {1} x _ {2}, x _ {1} x _ {3} + x _ {2} x _ {4}, x _ {3} ^ {2} - x _ {4} ^ {2}, x _ {3} x _ {4}, x _ {1} x _ {4} - x _ {2} x _ {3}\right)
\]

defines, on passing to the quotient, a homeomorphism of \(P_{3}\) into \(R^{6}\) .

\begin{enumerate}
\def\labelenumi{\arabic{enumi})}
\setcounter{enumi}{1}
\item
  Identify the projective space \(P_{n}\) with the quotient space of \(B_{n}\) defined in no. 1, Proposition 3, and let \(\varphi\) denote the canonical mapping of \(B_{n}\) onto \(P_{n}\) . A continuous mapping f of a topological space E into \(P_{n}\) is said to be inessential if there exists a continuous mapping g of E into \(B_{n}\) such that \(f = \varphi \circ g\) , and essential in the converse case (cf.~§ 2, Exercise 6). Show that there exists an entourage U of the uniformity of \(P_{n}\) such that, if f is an inessential mapping of a topological space E into \(P_{n}\) , every continuous mapping \(f'\) of E into \(P_{n}\) such that \((f(x), f'(x)) \in \mathbf{U}\) for all \(x \in E\) is also inessential.
\item
  If a continuous mapping of \(\mathbf{S}_1\) into \(\mathbf{P}_n\) is essential (Exercise 2) it cannot be extended to a continuous mapping of \(\mathbf{B}_2\) into \(\mathbf{P}_n\) (cf.~§ 2, Exercise 8).
\item
  If \(n > \mathbf{1}\) , there exists an essential mapping \(f\) of \(\mathbf{S}_1\) into \(\mathbf{P}_n\) such that \(f(\mathbf{S}_1) \neq \mathbf{P}_n\) {[}take \(f(\mathbf{S}_1)\) to be the image under \(\varphi\) of a diameter of \(\mathbf{B}_n\) ; this image is a projective line{]}. Deduce that, if \(n > \mathbf{1}\) , \(\mathbf{P}_n\) is homeomorphic to neither \(\mathbf{S}_n\) nor \(\mathbf{B}_n\) (use Exercise 3 of § 3 and Exercise 3 of § 2).
\item
  If we consider \(\mathbf{R}^2\) as embedded in \(\tilde{\mathbf{R}} \times \tilde{\mathbf{R}}\) , and \(\mathbf{R}\) as embedded in \(\tilde{\mathbf{R}}\) , show that the mapping \((x, y) \to x + y\) of \(\mathbf{R}^2\) into \(\mathbf{R}\) can be extended by continuity to the points \((\infty, a)\) and \((a, \infty)\) of \(\tilde{\mathbf{R}} \times \tilde{\mathbf{R}}\) for all finite values of \(a\) , but that it cannot be extended by continuity to \((\infty, \infty)\) . If we consider \(\mathbf{R}^2\) as embedded in \(\mathbf{P}_2\) ( \(\mathbf{R}\) still considered as embedded in \(\tilde{\mathbf{R}}\) ), \(x + y\) can be extended by continuity to all points of the line at infinity except the point with homogeneous coordinates \((1, -1, 0)\) .
\end{enumerate}

State and prove the analogues of these results for the product xy.

\begin{enumerate}
\def\labelenumi{\arabic{enumi})}
\setcounter{enumi}{5}
\item
  If \(n \geqslant 0\) , let \(\mathfrak{F}\) be the set of all closed subsets of \(\mathbf{P}_n\) endowed with the uniformity induced by that of \(\mathbf{P}_n\) by the procedure of Chapter II, § 2, Exercise 6 b). Show that the set \(\mathbf{P}_{n,p}\) of projective linear varieties of \(p \geqslant 0\) dimensions is closed in \(\mathfrak{F}\) , and that the topology induced on \(\mathbf{P}_{n,p}\) by the topology of \(\mathfrak{F}\) is the same as that defined in no. 5, Definition 2. {[}To define the entourages of the uniformity of \(\mathbf{P}_n\) we may consider \(\mathbf{P}_n\) as a quotient space of \(\mathbf{S}_n\) (no. 1, Proposition 2) and take finite coverings of \(\mathbf{S}_n\) by balls whose radii tend to o.{]}
\item
  \begin{enumerate}
  \def\labelenumii{\alph{enumii})}
  \tightlist
  \item
    Show that, in the notation of no. 5, the space \(\mathbf{L}_{n,p}\) is homeomorphic to the product of \(\mathbf{V}_{n,p}\) and \(\mathbb{R}^{p(p+1)/2}\) . Note that every matrix \(X\in\mathbf{L}_{n,p}\) can be expressed uniquely in the form \(U\) . \(\Upsilon\) , where \(\Upsilon\in\mathbf{V}_{n,p}\) and \(U=(u_{ij})\) is such that \(u_{ij}=0\) whenever \(i<j\) , and \(u_{ii}>0\) for \(\mathbf{i}\leqslant i\leqslant p\) . Put \(\Upsilon=f(X)\) .
  \end{enumerate}
\end{enumerate}

\begin{enumerate}
\def\labelenumi{\alph{enumi})}
\setcounter{enumi}{1}
\tightlist
\item
  In the notation of no. 5, Proposition 8, show that \(f\) is a homeomorphism of \(\mathbf{B}_{\sigma}\) onto \(f(\mathbf{B}_{\sigma})\) {[}note that \(X \to f(X_{\sigma}^{-1}X)\) is a continuous mapping of \(\mathbf{A}_{\sigma}\) onto \(f(\mathbf{B}_{\sigma})\) and that \(f(X_{\sigma}^{-1}X)\) belongs to the same class as \(X \mod \Delta_{n,p}\) ; then apply Lemma 2{]}. Deduce that the intersection \(\mathbf{D}_{\sigma}\) of \(\mathbf{A}_{\sigma}\) and \(\mathbf{V}_{n,p}\) is homeomorphic to the product of \(\mathbf{V}_{p,p}\) and \(\mathbb{R}^{p(n - p)}\) {[}every matrix belonging to \(\mathbf{D}_{\sigma}\) is uniquely expressible as the product of a matrix of \(\mathbf{V}_{p,p}\) and a matrix of \(f(\mathbf{B}_{\sigma})\) {]}.
\end{enumerate}

\begin{enumerate}
\def\labelenumi{\arabic{enumi})}
\setcounter{enumi}{7}
\tightlist
\item
  Let \(g\) be the mapping such that, for each matrix \(X \in \mathbf{L}_{n,p}\) , \(X'g(X)\) is the matrix formed by the first \(q\) rows of \(X\) ( \(q < p\) ). Show that the restriction of \(g\) to \(\mathbf{V}_{n,p}\) is a continuous open mapping of \(\mathbf{V}_{n,p}\) onto \(\mathbf{V}_{n,q}\) {[}use Exercise 7 a), by noting that if \(X = U.Y\) then
\end{enumerate}

\[
g (X) = U ^ {\prime}. g (\Upsilon),
\]

where \(U'\) is the matrix obtained from U by removing the rows and columns with indices \textgreater q; on the other hand, observe that f is open{]}.

\begin{enumerate}
\def\labelenumi{\arabic{enumi})}
\setcounter{enumi}{8}
\item
  Show that \(\mathbf{V}_{n,p}\) is connected if \(p < n\) (use Exercise 8 and Chapter I, § 11, no. 3, Proposition 7).
\item
  In a projective space \(P_{n}\) , let \(H_{n,p}\) be the ``quadric'' defined by the equation
\end{enumerate}

\[
x _ {0} ^ {2} + x _ {1} ^ {2} + \dots + x _ {p - 1} ^ {2} - x _ {p} ^ {2} - \dots - x _ {n} ^ {2} = 0 \quad (\mathrm{I} \leqslant p \leqslant n).
\]

Show that \(H_{n,1}\) and \(H_{n,n}\) are homeomorphic to \(S_{n-1}\) . If \(2 \leqslant p \leqslant n - 1\) , \(H_{n,p}\) is homeomorphic to the space obtained by identifying every point of the product \(S_{p-1} \times S_{n-p}\) (considered as a subspace of \(R^{p} \times R^{n-p+1}\) ) with its diametrically opposite point. Every point of \(H_{n,p}\) has an open neighbourhood homeomorphic to \(R^{n-1}\) .

Show that \(H_{3,2}\) is homeomorphic to \(S_{1} \times S_{1}\) {[}identify \(S_{1} \times S_{1}\) with \(T \times T\) , a pair of diametrically opposite points of \(S_{1} \times S_{1}\) being identified with a pair of points \((u, v)\) and \((1/2 + u, 1/2 + v)\) of \(T \times T\) ; then consider the mapping \((u, v) \to (u + v, u - v)\) of \(T \times T\) into itself{]}.

\begin{enumerate}
\def\labelenumi{\arabic{enumi})}
\setcounter{enumi}{10}
\tightlist
\item
  In a projective space \(P_{n}\) , let \(C_{n,p}\) be the ``quadric cone'' defined by the equation
\end{enumerate}

\[
x _ {1} ^ {2} + x _ {2} ^ {2} + \dots + x _ {p} ^ {2} - x _ {p + 1} ^ {2} - \dots - x _ {n} ^ {2} = 0 \quad (\mathrm{I} \leqslant p \leqslant n - \mathrm{I}).
\]

Show that the complement of \(\{0\}\) in \(C_{n,p}\) is homeomorphic to \(\mathbf{R} \times H_{n-1,p}\) , with the notation of Exercise 10.

